\documentclass{article}
\usepackage[spanish]{babel}
\usepackage[intoc, spanish]{nomencl}
\makenomenclature

\begin{document}
	\tableofcontents
	
	\section{Primera Sección}
	
	Lorem ipsum dolor sit amet, consectetuer adipiscing elit. Etiam lobortis facilisis sem. 
	Nullam nec mi et neque pharetra sollicitudin. Praesent imperdiet...
	
	% Define nomenclature entries
	\nomenclature{$c$}{Velocidad de la luz en el vacío}
	\nomenclature{$h$}{Constante de Planck}
	\nomenclature{$e$}{Carga elemental}
	\nomenclature{$m_e$}{Masa del electrón}
	\nomenclature{$\pi$}{Número pi (3.14159...)}
	\nomenclature{$G$}{Constante de gravitación universal}
	
	\section{Segunda Sección}
	
	La velocidad de la luz ($c$) es una constante fundamental en física. 
	La constante de Planck ($h$) es igualmente importante en mecánica cuántica.
	
	\nomenclature{$\alpha$}{Constante de estructura fina}
	\nomenclature{$\mu_0$}{Permeabilidad magnética del vacío}
	\nomenclature{$\epsilon_0$}{Permitividad eléctrica del vacío}
	
	\clearpage
	
	% Print the nomenclature
	\printnomenclature
	
\end{document}